%%%%%%%%%%  scan idx 320  |  folio 307  %%%%%%%%%%
\SGAPrintPage{307}%
\uline{is an isomorphism}.

To see that it is injective, we use the sequence of canonical homomorphisms:
\[
  \Hom(\Ghat,H\circ g)\longrightarrow
  \Hom(\Ghat\circ i,H)\longrightarrow
  \Hom(\Gg,H\circ g)\longrightarrow
  \Hom(\Ghat,H\circ g),
\]
the last arrow coming from the fact that \(\Ghat\) is a subfunctor of \(\Gg\). To verify that the composite homomorphism is the identity, we observe that a morphism of functors: \(\Ghat\longrightarrow H\circ g\) is determined by its value on the objects of \(\mathcal C\) whose filtration is coarse; indeed, let \(A\in\operatorname{Ob}(\mathcal C)\), and let \(A_g\) be the algebra \(A\) endowed with the coarse filtration; the canonical morphism \(A\to A_g\) defines a commutative diagram:
\[
\begin{tikzcd}
\Ghat(A) \arrow[r,"\Phi_A"] \arrow[d] &
H\circ g(A) \arrow[d]\\
\Ghat(A_g) \arrow[r,"\Phi_{A_g}"] &
H\circ g(A_g),
\end{tikzcd}
\]
showing that \(\Phi_A\) is the restriction of \(\Phi_{A_g}\) to \(\Ghat(A)\).

We next consider the composite homomorphism:
\[
  \Hom(\Ghat\circ i,H)\longrightarrow
  \Hom(\Gg,H\circ g)\longrightarrow
  \Hom(\Ghat,H\circ g)\longrightarrow
  \Hom(\Ghat\circ i,H),
\]
and it is clear that this is the identity, giving the asserted result.

\sgaStableTarget{sga6.e05.thm.1-10}\textbf{Theorem 1.10.} \uline{Let} \(H\) \uline{be a functor on the category} \(\mathcal C_0\) \uline{of} \(K\)-\uline{algebras}, \uline{with values in the category of separated complete filtered abelian groups}, \uline{and commuting with invariants under a finite group}. \uline{Then the map from} \(\Hom_{\underline{\mathrm{Ab}}\,\underline{\mathrm{fil}}}(\Ghat\circ i,H)\) \uline{to} \(H(K[X])\) \uline{which}, \uline{to a morphism of functors}, \uline{assigns}
%%%%%%%%%%  scan idx 321  |  folio 308  %%%%%%%%%%
\SGAPrintPage{308}%
\uline{its value on the element} \(1+Xt\) \uline{of} \(\Ghat\circ i(K[X])\), \uline{is injective}; \uline{its image is the subgroup of} \(H(K[X])\) \uline{formed by the elements} \(\theta\) \uline{such that} (\uline{with notation analogous to that of} 1.4 \uline{and} 1.5):
\[
\begin{array}{ll}
\text{i)}& H(f_1)(\theta)=0;\\[0.3em]
\text{ii)}& \forall n\geq1,\quad
H(f_n)\bigl(u_H^n(\theta)\bigr)\equiv u_H^n(\theta)
\ \mathrm{mod.}\ F^nH(K[X_1,\ldots,X_n]).
\end{array}
\]

The theorem follows immediately from Theorem 1.6 and Lemma 1.9, since \(H\circ g(K[X])\) does not depend on the filtration chosen on \(K[X]\).

The analogous theorem obtained by replacing \(\Ghat\) by \((\Ghat)^k\) is also true (cf. 1.7)

When there is no risk of confusion, we shall denote the functor \(\Ghat\circ i\) by \(\Ghat\); this is legitimate, since \(i\) identifies \(\mathcal C_0\) with a subcategory of \(\mathcal C\).

\sgaStableTarget{sga6.e05.para.1-11}1.11. We denote by \(\mathcal G\) the category of graded \(K\)-algebras, and by \(j\) the functor from \(\mathcal G\) to \(\mathcal C\) that assigns to a graded \(K\)-algebra \(B\) the \(K\)-algebra \(B\) endowed with the filtration canonically defined by its grading; the ``associated graded'' functor: \(\mathcal C\to\mathcal G\) will be denoted by \(\operatorname{Gr}\). There is an isomorphism of functors \(\varepsilon:\operatorname{Id}_{\mathcal G}\xrightarrow{\sim}\operatorname{Gr}\circ j\).

Let \(B\in\operatorname{Ob}(\mathcal G)\); we denote by \(1+\widehat B^+\) the multiplicative group of elements of \(\prod\limits_{i=0}^{\infty}B^i\) with augmentation \(1\). Endowing it with the filtration analogous to that of \(\Ghat\), we thus define a functor: \(\mathcal G\to\underline{\mathrm{Ab}}\,\underline{\mathrm{fil}}\).

Furthermore, for every \(A\in\operatorname{Ob}(\mathcal C)\), we define a homomorphism of filtered groups: \(\Ghat(A)\longrightarrow 1+\widehat{\operatorname{Gr}A}^{+}\), by sending the series \(1+\ldots+a_nt^n+\ldots\),
%%%%%%%%%%  scan idx 322  |  folio 309  %%%%%%%%%%
\SGAPrintPage{309}%
to the series \(1+\ldots+\bar a_n+\ldots\), where \(\bar a_n\) is the class of \(a_n\) in \(\operatorname{Gr}^nA\). Its kernel is the group \(\Ghat'(A)\) of formal series such that:
\[
  a_0=1;\qquad \forall i\geq1,\quad a_i\in\Fil^{i+1}A;
\]
we therefore have an isomorphism, functorial in \(A\):
\[
  \Ghat(A)/\Ghat'(A)\xrightarrow{\sim}1+\widehat{\operatorname{Gr}A}^{+}.
\]

\sgaStableTarget{sga6.e05.lem.1-12}\textbf{Lemma 1.12.} \uline{Let} \(H:\mathcal G\to\underline{\mathrm{Ab}}\,\underline{\mathrm{fil}}\) \uline{be a functor on} \(\mathcal G\); \uline{then the canonical homomorphism}:
\[
  \rho:\Hom_{\underline{\mathrm{Ab}}\,\underline{\mathrm{fil}}}(1+\widehat{\,\mathord{.}\,}^{+},H)
  \longrightarrow
  \Hom_{\underline{\mathrm{Ab}}\,\underline{\mathrm{fil}}}(\Ghat,H\circ\operatorname{Gr}),
\]
\uline{is an isomorphism}.

The homomorphism \(\rho\) is defined as the composite:
\[
  \Hom(1+\widehat{\,\mathord{.}\,}^{+},H)
  \longrightarrow
  \Hom(1+\widehat{\operatorname{Gr}\,\mathord{.}\,}^{+},H\circ\operatorname{Gr})
  \longrightarrow
  \Hom(\Ghat,H\circ\operatorname{Gr}),
\]
the last arrow being defined by the homomorphism \(\Ghat\longrightarrow1+\widehat{\operatorname{Gr}\,\mathord{.}\,}^{+}\)

We shall construct a morphism in the opposite direction:
\[
  \rho':\Hom(\Ghat,H\circ\operatorname{Gr})
  \longrightarrow
  \Hom(1+\widehat{\,\mathord{.}\,}^{+},H).
\]
For this, we shall show that every morphism of functors from \(\Ghat\) to \(H\circ\operatorname{Gr}\) factors through the canonical morphism \(\Ghat\longrightarrow1+\widehat{\operatorname{Gr}\,\mathord{.}\,}^{+}\), which, by 1.11, is equivalent to showing that it vanishes on \(\Ghat'\).

Let \(A\in\operatorname{Ob}(\mathcal C)\), \(\chi\in\Ghat'(A)\); we define a homomorphism of
%%%%%%%%%%  scan idx 323  |  folio 310  %%%%%%%%%%
\SGAPrintPage{310}%
filtered \(K\)-algebras \(K[X_i]_{i\geq1}\longrightarrow A\) by sending \(X_1\) to \(0\) and \(X_i\) to the \((i-1)\)-th coefficient of \(\chi\) for \(i>1\); using this homomorphism, we see that it suffices to verify that the morphism of functors under consideration vanishes on the element \(1+X_2t+\ldots+X_{n+1}t^n+\ldots\). We then use the endomorphism \(f\) of \(K[X_i]_{i\geq1}\) defined by \(f(X_i)=X_{i+1}\) for every \(i\), and the commutative diagram:
\[
\begin{tikzcd}
\Ghat(K[(X_i)]) \arrow[r] \arrow[d,"\Ghat(f)"'] &
H\circ\operatorname{Gr}(K[(X_i)]) \arrow[d,"H(\operatorname{Gr}(f))"]\\
\Ghat(K[(X_i)]) \arrow[r] &
H\circ\operatorname{Gr}(K[(X_i)]);
\end{tikzcd}
\]
now \(\operatorname{Gr}(f)=\operatorname{Gr}(f_1)\), where \(f_1\) is the augmentation homomorphism of \(K[X_i]_{i\geq1}\); since the latter annihilates \(\varphi\), the analogous diagram written for \(f_1\) shows that \(H(\operatorname{Gr}(f))\) annihilates the image of \(\varphi\) in \(H\circ\operatorname{Gr}(K[X_i]_{i\geq1})\), giving the result.

Passing to the quotient, we therefore associate to every morphism of functors: \(\Ghat\longrightarrow H\circ\operatorname{Gr}\) a morphism of functors: \(1+\widehat{\operatorname{Gr}\,\mathord{.}\,}^{+}\longrightarrow H\circ\operatorname{Gr}\); by composing with \(j\), and using the isomorphism \(\varepsilon\), we obtain a morphism: \(1+\widehat{\,\mathord{.}\,}^{+}\longrightarrow H\), which defines the required homomorphism \(\rho'\). It is immediate to verify that \(\rho\) and \(\rho'\) are inverses of one another.

\sgaStableTarget{sga6.e05.thm.1-13}\textbf{Theorem 1.13.} \uline{Let} \(H\) \uline{be a functor on} \(\mathcal G\), \uline{with values in the category of separated complete filtered abelian groups}, \uline{and commuting with invariants under a finite group}. \uline{Then the map from} \(\Hom_{\underline{\mathrm{Ab}}\,\underline{\mathrm{fil}}}(1+\widehat{\,\mathord{.}\,}^{+},H)\) \uline{to} \(H(K[X])\) \uline{which}, \uline{to a morphism of functors}, \uline{assigns its value on the element}
%%%%%%%%%%  scan idx 324  |  folio 311  %%%%%%%%%%
\SGAPrintPage{311}%
\(1+X\) of \(1+\widehat{K[X]}^+\) \uline{is injective}; \uline{its image is the subgroup of} \(H(K[X])\) \uline{formed by the elements} \(\theta\) \uline{such that} (\uline{notation of} 1.4,1.5):
\[
\begin{array}{ll}
\text{i)}& H(f_1)(\theta)=0;\\[0.3em]
\text{ii)}& \forall n\geq1,\quad
H(f_n)\bigl(u_H^n(\theta)\bigr)\equiv u_H^n(\theta)
\ \mathrm{mod.}\ F^nH(K[X_1,\ldots,X_n]).
\end{array}
\]

Immediate from 1.6 and 1.12.

There is an analogous statement on replacing \(\Ghat\) by \((\Ghat)^k\).

\sgaStableTarget{sga6.e05.rem.1-14}\textbf{Remark 1.14.} If \(H\) is a functor on \(\mathcal G\), with values in the category of abelian groups, commuting with invariants under a finite group, and satisfying the condition:
\[
  \forall B\in\operatorname{Ob}(\mathcal G),\qquad
  H(B)\xrightarrow{\sim}\mathop{\underleftarrow{\mathrm{Lim}}}\limits_n H(B/B^{\geq n})\quad(\text{where }B^{\geq n}=\prod_{i\geq n}B^i),
\]
then condition ii) of 1.13 is always satisfied when we take as the filtration on \(H\) the filtration by the kernels of the \(H(B)\longrightarrow H(B/B^{\geq n})\) (which is separated and complete). Indeed, let \(p_n\) be the canonical homomorphism
\[
  K[X_1,\ldots,X_n]\longrightarrow K[X_1,\ldots,X_n]/K[X_1,\ldots,X_n]^{\geq n};
\]
we have: \(p_n\circ f_n=p_n\), hence: \(H(p_n)\bigl[H(f_n)\circ u_H^n(\theta)-u_H^n(\theta)\bigr]=0\), that is:
\[
  H(f_n)\circ u_H^n(\theta)-u_H^n(\theta)
  \in F^nH(K[X_1,\ldots,X_n]).
\]

Thus, to specify a morphism of functors from \(1+\widehat{\,\mathord{.}\,}^{+}\) to \(H\), it suffices to specify an element of \(H(K[X])\) satisfying only condition i) of 1.13.

\sgaStableTarget{sga6.e05.para.1-15}1.15. Let \(H:\mathcal C\to\underline{\mathrm{Ab}}\,\underline{\mathrm{fil}}\) (resp. \(\mathcal C_0\to\underline{\mathrm{Ab}}\,\underline{\mathrm{fil}}\), \(\mathcal G\to\underline{\mathrm{Ab}}\,\underline{\mathrm{fil}}\)) be a functor satisfying the conditions of 1.3 (resp. 1.10, 1.13), and let \(\Phi:(\Ghat)^k\to H\) (resp. \((\Ghat\circ i)^k\to H\), \((1+\widehat{\,\mathord{.}\,}^{+})^k\to H\)) be a morphism of functors. We have seen
